% !TeX program = lualatex
\documentclass[UTF8,fontset=fandol,12pt,oneside]{ctexbook}
\usepackage[a4paper,left=20mm,right=20mm,top=24mm,bottom=22mm]{geometry}
\usepackage{amsmath}
\usepackage{liquidglass}
\usepackage[hidelinks]{hyperref}
\glasssettheoremlanguage{zh}
\definecolor{DemoPaper}{HTML}{F5F5F5}
\pagecolor{DemoPaper}

% The PDF can be rebuilt with the template's font pairing when the user's
% original files are available. The portable fallback is from TeX Live.
\providecommand\LiquidGlassFontRoot{fonts/}
\IfFileExists{\LiquidGlassFontRoot cnfonts/FZsongti_GBK.TTF}{%
  \setCJKmainfont[Path=\LiquidGlassFontRoot cnfonts/,
    BoldFont=FZheiti_GBK.TTF,ItalicFont=FZkaiti_GBK.TTF,
    SlantedFont=FZfangsong_GBK.TTF]{FZsongti_GBK.TTF}%
}{%
  \setCJKmainfont[BoldFont=FandolHei-Regular.otf,ItalicFont=FandolKai-Regular.otf,
    BoldItalicFont=FandolHei-Regular.otf]{FandolSong-Regular.otf}%
}
\IfFileExists{\LiquidGlassFontRoot enfonts/xits/XITS-Regular.otf}{%
  \setmainfont[Path=\LiquidGlassFontRoot enfonts/xits/,Extension=.otf,
    UprightFont=*-Regular,BoldFont=*-Bold,ItalicFont=*-Italic,
    BoldItalicFont=*-BoldItalic,Ligatures=TeX]{XITS}%
}{}
\linespread{1.25}
\pagestyle{empty}

\begin{document}
\setcounter{chapter}{3}
\setcounter{section}{4}
\noindent
\begin{lgproposition}{方法总结}{method}
   针对不同的问题背景,共用一套建模思路. ``五问建模法''
   \begin{lgtimeline}
      \lgtimelineitem 1. \textbf{问量：题目中究竟有哪些数量？}
   
   例如人数、件数、钱数、路程、时间、质量。要明确数量的单位、所属对象与时刻。
      \lgtimelineitem 2. \textbf{问关系：这些量如何联系？}
   
   是部分与整体、前后变化、差倍比较，还是利用数量公式进行换算？
      \lgtimelineitem 3. \textbf{问表达：怎样用未知数表示其他相关数量？}
   
   这里重点使用表格、线段图、状态变化图，而不是直接猜方程。
      \lgtimelineitem 4. \textbf{问约束：哪个条件必须成立？}
   
   这个条件是在限定总量、比较两个量，还是说明某个数量保持不变？
      \lgtimelineitem 5. \textbf{问方程：把有关表达式代入约束后得到什么？}
   
   检查等号两侧是否表示可比较的量，单位、对象、时刻和基准是否一致。
   \end{lgtimeline}
\end{lgproposition}

\clearpage
\section*{定理、定义、证明与推论}
\begin{lgdefinition}{算术平均数}{mean}
实数 $x_1,\ldots,x_n$ 的算术平均数为
$\bar{x}=\frac1n\sum_{k=1}^n x_k$。
\end{lgdefinition}

\begin{lgtheorem}{平方和恒等式}{squares}
\[\sum_{k=1}^n(x_k-\bar{x})^2=\sum_{k=1}^n x_k^2-n\bar{x}^{\,2}.\]
\end{lgtheorem}

\begin{lgproof}
展开平方，再使用定义~\ref{lgdef:mean} 代入 $\sum_k x_k=n\bar{x}$。
\end{lgproof}

\begin{lgcorollary}{方差非负}{variance}
由定理~\ref{lgthm:squares}，方差为零当且仅当各个数相等。
\end{lgcorollary}

\clearpage
\section*{其他定理种类}
\begin{lglemma}{两数的平方和}{two}
$(a-b)^2\ge0$ 给出 $a^2+b^2\ge 2ab$。
\end{lglemma}
\begin{lgexample}{代入检查}{check}
$a=3$、$b=1$ 时，左侧为 $10$，右侧为 $6$。
\end{lgexample}
\begin{lgexercise}{等号条件}{equality}
由引理~\ref{lglem:two} 找出等号成立的条件。
\end{lgexercise}
\begin{lgproblem}{最小值}{minimum}
若 $a+b=6$，求 $a^2+b^2$ 的最小值。
\end{lgproblem}
\begin{lgsolution}
用 $(a-b)^2=2(a^2+b^2)-(a+b)^2$，得最小值 $18$。
\end{lgsolution}
\begin{lgremark*}{应用约束}
把数量代入等式前，先核对单位。
\end{lgremark*}
\begin{lgapplication}{数量比较}{comparison}
使用平均数时，明确参与平均的样本。
\end{lgapplication}
\begin{lgconjecture}[lg theme=dark]{深色主题示例}{dark}
观察更多数据，再决定猜想是否值得继续验证。
\end{lgconjecture}
\end{document}
